\documentclass[12pt]{article}

\usepackage[margin=1in]{geometry}
\usepackage{setspace}
\usepackage{parskip}
\usepackage{hyperref}
\usepackage{enumitem}
\usepackage{titlesec}

\titleformat{\section}{\large\bfseries}{}{0em}{}[\titlerule]
\titlespacing{\section}{0pt}{1.2em}{0.6em}

\onehalfspacing

\begin{document}

\begin{center}
    {\LARGE \textbf{Research Statement}} \\[0.8em]
    {\large [Bohdan Leniuk]} \\[0.3em]
    {\normalsize [University of Georgia] \quad $\cdot$ \quad [bohdan.leniuk@uga.edu]}
\end{center}

\vspace{1.5em}

My research lies at the intersection of macroeconomics and labor economics,
with a focus on job search, worker mobility, sorting, and compensation.
I am interested in how labor market frictions and institutional features---both
informational and contractual---shape where workers sort across employers and
how they are paid, and what the resulting implications are for efficiency and
inequality.
My current work pursues these questions through two complementary projects:
a structural study of how promotions and unemployment shape sorting under
asymmetric information, and a reduced-form empirical investigation of how
retirement plan design affects worker mobility and wages.

%------------------------------------------------------------------
\section{Job Market Paper: Promotions, Unemployment, and Labor Market Sorting Under Asymmetric Information}
%------------------------------------------------------------------

My job market paper studies how promotions and unemployment jointly determine
labor market sorting and wages.
Using the NLSY79, I document that previously promoted workers face
disproportionately large wage losses upon re-employment---averaging 15\%
below pre-separation wages versus 1\% for non-promoted workers---and that
these losses grow significantly with unemployment duration, at a rate of
0.55 percentage points per month compared to 0.14 for their non-promoted
counterparts.
I also find that a substantial share of promotions are awarded to workers who
are plausibly low-ability, suggesting that the promotion signal itself is noisy
and that firms face a non-trivial assignment problem both internally and when
hiring from outside.
I interpret these facts through a structural model in which firms learn about
worker ability from employment history.
Promotions serve as positive signals of ability, but prolonged unemployment
acts as a noisy counteracting signal, causing many formerly promoted workers
to be re-employed in lower-level positions---a form of downward mobility
driven entirely by information frictions rather than changes in productivity.
The model also delivers policy implications for unemployment insurance design:
benefit profiles tied to past wages worsen sorting and reduce aggregate output
by distorting the informativeness of unemployment duration as an ability
signal, while policies targeted toward low-ability workers improve sorting.

%------------------------------------------------------------------
\section{Work in Progress: Retirement Plan Design, Worker Mobility, and Wages}
%------------------------------------------------------------------

\textit{Joint with Martin Gervais.}

\medskip

This project examines how employer-sponsored retirement plan design affects
worker mobility and wages.
We construct a novel dataset by scraping retirement plan information for all
U.S.\ firms with more than 1{,}000 employees, capturing detailed variation
in vesting schedule structures and employer contribution generosity.
To our knowledge, no existing dataset provides this level of coverage of
plan-level design features across the near-universe of large U.S.\ employers.
Linking this to employer--employee matched administrative data, we pursue a
reduced-form strategy that exploits cross-firm variation in plan design to
identify how vesting schedules and benefit generosity affect turnover and wages.
Vesting provisions create explicit financial penalties for early separation
that may suppress voluntary job mobility, particularly around vesting cliffs,
while more generous retirement contributions may be partially offset by lower
wages---a form of compensating differential operating through the non-wage
margin.
By documenting these relationships, the paper contributes to our understanding
of how non-wage compensation shapes equilibrium labor market outcomes.

%------------------------------------------------------------------
\section{Future Research}
%------------------------------------------------------------------

Because detailed retirement plan data of this kind is essentially absent from
the existing literature, the dataset assembled above opens many doors for
original and valuable research.
Among the most promising directions are questions about the frictions and
inefficiencies that current plan designs introduce into the labor market.
For instance, vesting schedules may distort workers' optimal separation
decisions, generating inefficient lock-in that a structural model could
quantify and evaluate against alternative plan designs.
Separately, contribution formulas tied to salary and lengthy vesting
requirements raise questions about the distributional consequences of plan
design, as these features may disproportionately benefit higher earners and
compound existing wage inequality through the retirement savings channel.
Beyond these two directions, the data is rich enough to speak to questions
about firm-level sorting, the role of benefits in job search, and the
interaction between retirement plan generosity and wage bargaining.
These represent just a subset of the questions this dataset is well-positioned
to address.

\end{document}
